% answers_short.tex —— 2010 年·普通高等学校招生考试·辽宁卷·文科 · 简短答案速查
% 由 tools/make_short_answers.py 自 answers.tex 的 \gksolhead/\gkanswer 提取，
% 并与 2010/2010_数学真题_辽宁文_答案要点.md 逐题比对后生成，勿手改（改源 MD 后重跑）。

\begin{gknotice}
\textbf{使用说明}：本节只给最终答案便于快速核对，完整推导见解析版。答案为编写者独立推导，\textbf{非官方标准答案}。
\end{gknotice}

\gkpart{选择题}
\begin{center}\begin{tabular}{ccccccccccccc}
\toprule
题号 & 1 & 2 & 3 & 4 & 5 & 6 & 7 & 8 & 9 & 10 & 11 & 12 \\
\midrule
答案 & D & A & B & C & B & C & B & C & D & A & A & D \\
\bottomrule
\end{tabular}\end{center}

\gkpart{填空题}
\begin{itemize}[leftmargin=2.6em]
  \item[\textbf{13.}] $\dfrac13$。
  \item[\textbf{14.}] $15$。
  \item[\textbf{15.}] $(3,8)$。
  \item[\textbf{16.}] $2\sqrt3$。
\end{itemize}

\gkpart{解答题}
\begin{itemize}[leftmargin=2.6em]
  \item[\textbf{17.}] （1）$A=\dfrac{2\pi}{3}$；（2）$\triangle ABC$ 为等腰三角形。
  \item[\textbf{18.}] （1）频率/组距依次为 $A$：$0.06,0.08,0.04,0.02$，$B$：$0.02,0.05,0.04,0.06,0.03$；中位数 $A$ 为 $67.5$、$B$ 为 $73.75$，故 $A$ 的中位数较小；（2）$a=70$，$b=30$，$c=35$，$d=65$，$n=200$，有 $99.9\%$ 的把握认为有差异。
  \item[\textbf{19.}] （1）证明见解析；（2）$A_1D:DC_1=1$。
  \item[\textbf{20.}] （1）焦距为 $4$；（2）$\dfrac{x^2}{9}+\dfrac{y^2}{5}=1$。
  \item[\textbf{21.}] （1）见解析；（2）证明见解析。
\end{itemize}

\gkpart{选考题}
\begin{itemize}[leftmargin=2.6em]
  \item[\textbf{22.}] （1）证明见解析；（2）$\angle BAC=90^\circ$。
  \item[\textbf{23.}] （1）$M\left(\dfrac{\pi}{3},\dfrac{\pi}{3}\right)$；（2）$\begin{cases}x=1+\left(\dfrac{\pi}{6}-1\right)t,\\y=\dfrac{\sqrt3\pi}{6}t,\end{cases}$（$t$ 为参数）。
  \item[\textbf{24.}] 证明见解析；当且仅当 $a=b=c=\sqrt[4]{3}$ 时等号成立。
\end{itemize}
